% Three Problems + Spec + Bugs + Unified + Pipeline (Ch 32-36)
% Lines 7204-7829 of main.tex
% This file is for agent editing — will be merged back

\chapter{The Three Problems, One Framework}
\label{ch:three-problems}

The cubical--cohomological framework solves three problems with the
\emph{same} mathematical machinery:

\begin{center}
\begin{tabular}{l l l}
  \textbf{Problem} & \textbf{CTTT formulation} & \textbf{Z3 query} \\
  \hline
  Equivalence & $\Path_{A \to B}(f, g)$ exists? &
    $\exists x.\, \den{f}(x) \neq \den{g}(x)$? \\
  Spec checking & $\Path_{A \to B}(f, \mathsf{spec})$ exists? &
    $\exists x.\, \lnot\,S(x, \den{f}(x))$? \\
  Bug finding & $\Hone(\covers, \Iso(\Sem_f, \Sem_\mathsf{spec})) \neq 0$? &
    which fiber has the obstruction? \\
\end{tabular}
\end{center}

\textbf{Equivalence} asks: is there a cubical path from $f$ to $g$?

\textbf{Spec checking} asks: is there a cubical path from $f$ to
\emph{any} program satisfying the specification $S$?  Since $S$ is
deterministic (at most one valid output per input), this reduces to:
does $f$ itself satisfy $S$?

\textbf{Bug finding} asks: on which type fiber does $f$ \emph{fail}
to satisfy $S$?  The cohomological decomposition localizes the bug
to a specific fiber --- a specific type of input where the program
misbehaves.

All three use:
\begin{enumerate}
  \item The \textbf{PyObj Z3 theory} to compile programs to Z3 terms.
  \item The \textbf{type fibration} to decompose the input space.
  \item The \textbf{\v{C}ech complex} to relate local and global
    properties.
  \item The \textbf{path constructors} to bridge structural gaps.
\end{enumerate}

The only difference is what we compare the program's denotation
\emph{against}: another program (equivalence), a specification
(spec checking), or an ideal behavior (bug finding).


\chapter{Specification Checking}
\label{ch:spec-checking}

\section{Specifications as Dependent Types}

In CTTT, a specification is a \emph{dependent type over the value
presheaf}:

\begin{definition}[Specification]
A specification $S$ for a function type $A \to B$ is a type family:
\[
  S : \prod_{x : A} B \to \Type
\]
where $S(x, y)$ is inhabited iff $y$ is a valid output for input $x$.
\end{definition}

A program $f : A \to B$ \emph{satisfies} $S$ iff:
\[
  \vdash \lambda x.\, \mathsf{witness}(x) :
  \prod_{x : A} S(x, f(x))
\]
--- for every input $x$, there is a witness that $f(x)$ is valid.

In Z3, this becomes: is there any $x$ where $\lnot\,S(x, \den{f}(x))$?
\begin{itemize}
  \item \textbf{UNSAT}: $f$ satisfies $S$ for all inputs.
    The UNSAT certificate is the constructive witness.
  \item \textbf{SAT}: $f$ violates $S$ on the counterexample $x$
    from the Z3 model.  This is a \emph{bug}.
\end{itemize}

\section{Specifications from Natural Language}

Deppy's \texttt{@guarantee} decorator attaches a natural-language
specification to a function:
\begin{lstlisting}
@guarantee("returns a sorted list with no duplicates")
def unique_sorted(lst):
    return sorted(set(lst))
\end{lstlisting}

The NL specification is parsed into a conjunction of Z3 predicates:
\begin{align}
  S_1(x, y) &: \mathsf{is\_sorted}(y)
    && \text{``sorted''} \\
  S_2(x, y) &: \mathsf{no\_duplicates}(y)
    && \text{``no duplicates''} \\
  S_3(x, y) &: \mathsf{subset}(y, x)
    && \text{``from the input''} \\
  S(x, y) &= S_1 \land S_2 \land S_3
\end{align}

Each predicate is encoded using the shared function symbols:
\begin{itemize}
  \item $\mathsf{is\_sorted}(y) \equiv \mathsf{py\_sorted}(y) = y$
    (using the $\mathsf{sorted\_idem}$ axiom).
  \item $\mathsf{no\_duplicates}(y) \equiv \mathsf{py\_set}(y) = y$
    (using the $\mathsf{set\_idem}$ axiom).
  \item $\mathsf{subset}(y, x) \equiv \mathsf{py\_set}(y) \subseteq
    \mathsf{py\_set}(x)$.
\end{itemize}

\section{Fiber-Wise Spec Checking}

The cohomological decomposition applies to spec checking:

\begin{definition}[Specification Presheaf]
The \emph{specification presheaf} $\mathsf{Spec}_S$ assigns to each
type fiber $U$ the set of valid behaviors:
\[
  \mathsf{Spec}_S(U) = \{(x, y) \in U \times B \mid S(x, y)\}
\]
\end{definition}

The program $f$ satisfies $S$ on fiber $U$ iff:
\[
  \forall x \in U.\; (x, \den{f}(x)) \in \mathsf{Spec}_S(U)
\]

This is checked fiber-by-fiber, just like equivalence:
\begin{enumerate}
  \item For each fiber $U_\tau$, ask Z3:
    $\exists x.\, \mathsf{tag}(x) = \tau \land \lnot\,S(x, \den{f}(x))$.
  \item Collect local verdicts.
  \item Compute $\Hone$ to check if local satisfaction glues globally.
\end{enumerate}

The \v{C}ech decomposition is especially valuable for spec checking
because specifications often have \emph{different constraints on
different types}:
\begin{itemize}
  \item On the $\mathsf{Int}$ fiber: ``output is non-negative.''
  \item On the $\mathsf{Str}$ fiber: ``output is non-empty.''
  \item On the $\mathsf{List}$ fiber: ``output has the same length.''
\end{itemize}

Each constraint is checked independently on its fiber, then glued.

\section{Example: Sorting Specification}

\begin{lstlisting}[caption={Sorting with spec check}]
@guarantee("returns a sorted permutation of the input")
def my_sort(lst):
    result = []
    for x in lst:
        # BUG: inserts at wrong position
        result.insert(0, x)
    return result
\end{lstlisting}

\textbf{Compilation}: $\den{\mathsf{my\_sort}}(p_0)$ produces a
sequence built by repeated $\mathsf{py\_meth\_insert}(0, x)$ ---
which prepends each element, producing the \emph{reverse} of the input.

\textbf{Spec}: $S(x, y) = \mathsf{py\_sorted}(y) = y \;\land\;
\mathsf{len}(y) = \mathsf{len}(x)$.

\textbf{Z3 check}: Is there an $x$ where $\mathsf{py\_sorted}(\den{f}(x))
\neq \den{f}(x)$?  Z3 finds a counterexample: $x = [2, 1]$, where
$\den{f}([2,1]) = [1, 2]$ is actually sorted (by accident!), but
$x = [1, 2]$ gives $\den{f}([1,2]) = [2, 1]$ which is NOT sorted.
\textbf{SAT} --- bug found.

\textbf{Fiber localization}: The bug is on the $\mathsf{Pair}$ fiber
(list inputs).  $\Hone(\{\mathsf{Pair}\}, \Iso(\Sem_f, \Sem_S)) \neq 0$.

\section{Example: Fibonacci Specification}

\begin{lstlisting}[caption={Fibonacci with off-by-one bug}]
@guarantee("returns the n-th Fibonacci number")
def fib(n):
    if n <= 1: return n
    a, b = 0, 1
    for _ in range(n):  # BUG: should be range(n-1)
        a, b = b, a + b
    return a
\end{lstlisting}

\textbf{Spec}: $S(n, y)$ is the Fibonacci recurrence:
$F(0) = 0$, $F(1) = 1$, $F(n) = F(n-1) + F(n-2)$.

\textbf{Compilation}: The loop runs $n$ times instead of $n-1$,
producing $F(n+1)$ instead of $F(n)$ for $n \geq 2$.

\textbf{Z3 check}: $\exists n.\, \mathsf{tag}(n) = \mathsf{TInt}
\land n \geq 0 \land \den{f}(n) \neq F(n)$.  Z3 finds $n = 2$:
$\den{f}(2) = 2$ but $F(2) = 1$.  \textbf{SAT} --- bug found.

\textbf{Fiber localization}: Bug is on the $\mathsf{Int}$ fiber for
$n \geq 2$.  The base cases $n = 0, 1$ pass (the loop doesn't execute
or executes once correctly).

\section{Spec Checking via Equivalence}

A powerful technique: reduce spec checking to equivalence checking.

\begin{theorem}[Spec Checking as Equivalence]
\label{thm:spec-as-equiv}
Let $S$ be a deterministic specification with a known \emph{reference
implementation} $r$ (a program that satisfies $S$).  Then $f$ satisfies
$S$ iff $f \equiv r$.
\end{theorem}

\begin{proof}
Since $S$ is deterministic, there is exactly one valid output per
input.  Both $f$ and $r$ must produce this output.  So $f$ satisfies
$S$ iff $f(x) = r(x)$ for all $x$ iff $f \equiv r$.
\end{proof}

This reduces spec checking to the equivalence problem, which we
already solve with the full cubical--cohomological machinery
(path constructors, dichotomy theories, \v{C}ech complex).

\begin{example}
To check that a sorting implementation is correct, compare it against
\texttt{sorted} (the reference implementation).  The shared function
symbol $\mathsf{py\_sorted}$ and the sorted axioms handle the
comparison.
\end{example}


\chapter{Bug Finding via Cohomological Obstruction}
\label{ch:bug-finding}

\section{Bugs as $\Hone$ Obstructions}

A \emph{bug} is a violation of a specification.  In CTTT, bugs are
\emph{cohomological obstructions}: elements of $\Hone$ that prevent
local correctness from gluing to global correctness.

\begin{definition}[Bug as Obstruction]
A bug in program $f$ with respect to specification $S$ is a
non-trivial element of:
\[
  \Hone(\covers, \Iso(\Sem_f, \Sem_S))
\]
where $\Sem_S$ is the presheaf of the specification (the ``ideal''
behavior).
\end{definition}

The obstruction element tells us:
\begin{enumerate}
  \item \textbf{Which fiber}: the type of input that triggers the bug.
  \item \textbf{Which overlap}: whether the bug is a local failure
    (wrong answer on one fiber) or a gluing failure (correct on each
    fiber separately but inconsistent across fibers).
  \item \textbf{The counterexample}: the Z3 model gives a concrete
    input that exhibits the bug.
\end{enumerate}

\section{Bug Localization}

The \v{C}ech decomposition provides automatic \emph{bug localization}:

\begin{theorem}[Bug Localization via Fiber Decomposition]
If $\Hone(\covers, \Iso(\Sem_f, \Sem_S)) \neq 0$, the obstruction
is localized to specific fibers $U_{i_1}, \ldots, U_{i_k}$ where
the local checks fail.  The Z3 counterexamples on these fibers are
the \emph{minimal bug-triggering inputs}.
\end{theorem}

\begin{example}[Type-Dependent Bug]
\begin{lstlisting}
@guarantee("returns the absolute value")
def my_abs(x):
    if x < 0:
        return -x
    return x
\end{lstlisting}

This is correct on the $\mathsf{Int}$ fiber.  But on the
$\mathsf{Str}$ fiber, the comparison \texttt{x < 0} is meaningless
(Python 3 raises \texttt{TypeError} for \texttt{str < int}).

The \v{C}ech decomposition detects:
\begin{itemize}
  \item $\Iso(\mathsf{Int})$: UNSAT (correct on integers).
  \item $\Iso(\mathsf{Str})$: SAT (bug on strings --- comparison
    with integer fails).
  \item $\Hone \neq 0$ on the $\mathsf{Str}$ fiber.
\end{itemize}

The bug is localized to the string fiber.  The fix: add a type guard
\texttt{if not isinstance(x, (int, float)): raise TypeError(...)}.
\end{example}

\section{Bug Classification by Cohomological Degree}

Bugs can be classified by which cohomological degree they appear in:

\begin{center}
\begin{tabular}{l l l}
  \textbf{Degree} & \textbf{Meaning} & \textbf{Example} \\
  \hline
  $\Hzero = 0$ & No correct behavior at all &
    Function always crashes \\
  $\Hzero \neq 0, \Hone \neq 0$ &
    Correct on some fibers, wrong on others &
    Off-by-one for large $n$ \\
  $\Hone = 0, \Htwo \neq 0$ &
    Locally correct, inconsistent gluing &
    Type coercion bug \\
\end{tabular}
\end{center}

\textbf{$\Hzero = 0$ bugs}: The program fails the specification on
\emph{every} fiber.  This is a total failure --- the algorithm is
fundamentally wrong.

\textbf{$\Hzero \neq 0, \Hone \neq 0$ bugs}: The program is correct
on some fibers but not others.  This is the most common case ---
edge cases, boundary conditions, type-specific failures.

\textbf{$\Hone = 0, \Htwo \neq 0$ bugs}: These are subtle --- the
program is correct on each fiber individually, but the transition
functions between fibers are inconsistent.  This occurs when a
program handles type coercions incorrectly (e.g., treating
\texttt{True} as \texttt{1} in one branch but as a boolean in
another).

\section{The Bug-Finding Pipeline}

\begin{enumerate}
  \item \textbf{Input}: program $f$ + specification $S$ (NL or formal).
  \item \textbf{Compile}: $\den{f}$ and $\den{S}$ to Z3 terms.
  \item \textbf{Decompose}: build type site, identify fibers.
  \item \textbf{Local check}: for each fiber $U_\tau$, ask Z3:
    $\exists x \in U_\tau.\, \lnot\,S(x, \den{f}(x))$.
  \item \textbf{Report}:
    \begin{itemize}
      \item If all fibers pass: ``No bugs found (verified on all
        type fibers).''
      \item If fiber $U_\tau$ fails: ``Bug on $\tau$ inputs:
        $f(\mathsf{counterexample}) = \mathsf{wrong\_value}$,
        expected to satisfy $S$.''
    \end{itemize}
  \item \textbf{Localize}: rank fibers by severity (how many section
    pairs fail, how deep the obstruction).
\end{enumerate}

\section{Example: Off-by-One in Range}

\begin{lstlisting}[caption={Sum with off-by-one}]
@guarantee("returns the sum 1 + 2 + ... + n")
def sum_to_n(n):
    total = 0
    for i in range(n):  # BUG: should be range(1, n+1)
        total += i
    return total
\end{lstlisting}

\textbf{Spec}: $S(n, y) = (y = n \cdot (n+1) / 2)$.

\textbf{Compilation}: The loop computes $\fold(\mathsf{ADD}, 0, n,
\mathsf{IntObj}(0))$, which is $0 + 1 + \ldots + (n-1) =
n(n-1)/2$.

\textbf{Z3 check}: $\exists n.\, n \geq 0 \land n(n-1)/2 \neq
n(n+1)/2$.  Z3 finds $n = 2$: program returns $1$, spec says $3$.

\textbf{Obstruction}: $\Hone(\mathsf{Int}) \neq 0$.  The bug is
an off-by-one in the range bounds --- the fold starts at $0$
instead of $1$ and stops at $n$ instead of $n+1$.

\section{Example: Mutation Bug (False Aliasing)}

\begin{lstlisting}[caption={Mutable default argument}]
@guarantee("returns a list containing just x")
def make_list(x, lst=[]):
    lst.append(x)
    return lst
\end{lstlisting}

\textbf{Spec}: $S(x, y) = (y = [x])$.

\textbf{Bug}: The mutable default argument \texttt{lst=[]} is shared
across calls.  On the first call, $f(1) = [1]$ (correct).  On the
second call, $f(2) = [1, 2]$ (wrong --- the list persists).

\textbf{CTTT analysis}: The parameter \texttt{lst} has type
$\mathsf{Ref}$ (mutable).  On the $\mathsf{Ref}$ fiber, the
mutation $\mathsf{py\_mut\_append}$ modifies a shared heap location.
The spec says the output should be $[\mathsf{x}]$, but the Z3 term
shows accumulation: $\mathsf{mut\_append}(\mathsf{mut\_append}(\ldots,
x_1), x_2)$.

\textbf{Obstruction}: $\Hone(\mathsf{Ref}) \neq 0$.  The bug is on
the $\mathsf{Ref}$ fiber --- exactly where mutation effects matter.


\chapter{Unified Theory: Equivalence, Specs, and Bugs}
\label{ch:unified}

\section{The Verification Triangle}

The three problems form a triangle connected by the CTTT framework:

\begin{center}
\begin{tikzcd}[column sep=4em, row sep=3em]
  \text{Equivalence} \arrow[r, "\text{spec as ref.\ impl.}"]
    \arrow[d, "\text{bug = failed equiv.}"'] &
  \text{Spec Checking} \arrow[dl, "\text{bug = spec violation}"] \\
  \text{Bug Finding} &
\end{tikzcd}
\end{center}

\begin{itemize}
  \item \textbf{Equivalence $\to$ Spec Checking}: Given a reference
    implementation $r$ satisfying spec $S$, check $f \equiv r$
    (Theorem~\ref{thm:spec-as-equiv}).
  \item \textbf{Spec Checking $\to$ Bug Finding}: If the spec check
    fails, the counterexample is a bug.
  \item \textbf{Bug Finding $\to$ Equivalence}: If two programs have
    the \emph{same} bugs (same $\Hone$ obstructions on the same
    fibers), they might still be equivalent (``equivalently buggy'').
\end{itemize}

\section{The Universal Verification Query}

All three problems reduce to a single Z3 query form:
\[
  \exists x.\; \mathsf{tag}(x) \in F \;\land\;
  \den{f}(x) \neq \mathsf{target}(x)
\]

where:
\begin{itemize}
  \item $F$ is the set of allowed type fibers (from duck type
    inference).
  \item $\mathsf{target}$ is:
    \begin{itemize}
      \item $\den{g}$ for equivalence checking,
      \item the specification's unique output for spec checking,
      \item the ``ideal'' behavior for bug finding.
    \end{itemize}
\end{itemize}

The cohomological decomposition splits this into per-fiber queries.
The path constructors (equational axioms) help Z3 prove that
$\den{f}(x) = \mathsf{target}(x)$ even when the terms are
structurally different.

\section{Composition: Spec-Guided Equivalence}

A powerful composition: use a specification to \emph{guide}
equivalence checking.

\begin{theorem}[Spec-Guided Equivalence]
If both $f$ and $g$ satisfy a deterministic specification $S$,
then $f \equiv g$.
\end{theorem}

This is Theorem~\ref{thm:det-spec} from Chapter~\ref{ch:algo-deep},
but now we have the full pipeline to check it:
\begin{enumerate}
  \item Check $f \models S$ (spec checking via Z3).
  \item Check $g \models S$ (spec checking via Z3).
  \item If both pass, conclude $f \equiv g$ (by deterministic
    specification uniqueness).
\end{enumerate}

This handles the hardest equivalence cases (dichotomy D20: different
algorithms for the same problem) by reducing them to two independent
spec checks.

\section{Composition: Equivalence-Guided Bug Finding}

Another composition: use a known-correct program to find bugs in
a new implementation.

\begin{theorem}[Differential Testing via CTTT]
Given a reference implementation $r$ (known correct) and a new
implementation $f$, the equivalence check $f \equiv r$ either:
\begin{enumerate}
  \item Succeeds ($\Hone = 0$): $f$ is correct (no bugs).
  \item Fails ($\Hone \neq 0$): the obstruction localizes the bug
    to a specific fiber, with a concrete counterexample.
\end{enumerate}
\end{theorem}

This is \emph{differential testing without execution}: we compare
Z3 denotations, not runtime outputs.  The counterexample is a
\emph{symbolic} input, not a random sample.

\section{The Graded Verification Verdict}

The full CTTT verdict is richer than a simple pass/fail:

\begin{definition}[Graded Verification Verdict]
The verdict for program $f$ against target $T$ (another program or
a specification) is a triple:
\[
  (\Hzero, \Hone, \mathsf{obstructions})
\]
where:
\begin{itemize}
  \item $\Hzero$ = number of fibers where $f$ matches $T$
    (``how much is correct'').
  \item $\Hone$ = number of obstructions (``how many bugs'').
  \item $\mathsf{obstructions}$ = list of
    $(\mathsf{fiber}, \mathsf{counterexample})$ pairs
    (``where and what are the bugs'').
\end{itemize}
\end{definition}

A perfect program has $(\Hzero = |\covers|, \Hone = 0, [])$.
A totally wrong program has $(\Hzero = 0, \Hone = |\covers|,
[\ldots])$.
A program with one edge-case bug has
$(\Hzero = |\covers| - 1, \Hone = 1, [(\tau, x)])$.

This graded verdict is more informative than ``correct'' or
``incorrect'' --- it tells the developer \emph{exactly} where the
program fails and on what type of input.


% ══════════════════════════════════════════════════════════
\part{Implementation Architecture}
% ══════════════════════════════════════════════════════════

\chapter{The Verification Pipeline}
\label{ch:pipeline}

\section{Pipeline Overview}

\begin{enumerate}
  \item \textbf{Parse}: Python source $\to$ AST.
  \item \textbf{Compile}: AST $\to$ Z3 $\PyObj$ terms (sections).
  \item \textbf{Duck type inference}: narrow type fibers per parameter.
  \item \textbf{Build site category}: type fibers $\to$ sites +
    morphisms + overlaps.
  \item \textbf{Local equivalence}: per-fiber Z3 SAT check.
  \item \textbf{\v{C}ech $\Hone$}: full cohomology computation.
  \item \textbf{Verdict}: $\Hone = 0 \Rightarrow$ Equivalent;
    $\Hone \neq 0 \Rightarrow$ Inequivalent.
\end{enumerate}

\section{AST Compilation}

The compiler walks the Python AST depth-first, producing Z3 terms:

\begin{center}
\small
\begin{tabular}{l l}
  \textbf{Python construct} & \textbf{Z3 term} \\
  \hline
  Integer literal $n$ & $\mathsf{IntObj}(n)$ \\
  String literal $s$ & $\mathsf{StrObj}(s)$ \\
  Variable $x$ & $\mathsf{env}[x]$ (Z3 Const) \\
  $x + y$ & $\mathsf{binop}(\mathsf{ADD}, \den{x}, \den{y})$ \\
  $-x$ & $\mathsf{unop}(\mathsf{NEG}, \den{x})$ \\
  $x > y$ & $\mathsf{binop}(\mathsf{GT}, \den{x}, \den{y})$ \\
  $x$ \texttt{and} $y$ & $\mathsf{If}(\mathsf{truthy}(\den{x}), \den{y}, \den{x})$ \\
  $x$ \texttt{or} $y$ & $\mathsf{If}(\mathsf{truthy}(\den{x}), \den{x}, \den{y})$ \\
  \texttt{if c: t else: f} & $\mathsf{If}(\mathsf{truthy}(\den{c}), \den{t}, \den{f})$ \\
  \texttt{sorted(x)} & $\mathsf{py\_sorted}(\den{x})$ (shared) \\
  \texttt{x.append(v)} & $\mathsf{py\_mut\_append}(\den{x}, \den{v})$ (shared) \\
  \texttt{isinstance(x, int)} & $\mathsf{BoolObj}(\mathsf{tag}(\den{x}) = \mathsf{TInt})$ \\
  \texttt{for i in range(n)} & $\fold$ or RecFunction \\
  Recursive $f(n-1)$ & RecFunction or canonical fold \\
  \texttt{[e for x in xs]} & $\mathsf{py\_comp}_h(\den{xs})$ (shared) \\
\end{tabular}
\end{center}

\section{Section Extraction}

The compiler produces a list of \emph{sections} --- pairs
$(\mathsf{guard}, \mathsf{term})$:

\begin{itemize}
  \item Each \texttt{return} statement in the function body produces
    one section.
  \item The \texttt{guard} is the conjunction of all branch conditions
    leading to that return.
  \item The \texttt{term} is the Z3 expression for the return value.
\end{itemize}

These sections form the value presheaf: each section is a
``local behavior'' of the program under a specific set of conditions.

\section{Site Category Construction}

The type site is built from duck type inference:
\begin{enumerate}
  \item For each parameter, analyze its usage in both programs to
    determine allowed types (int, str, list, etc.).
  \item Form the Cartesian product of allowed types across parameters.
  \item Each element of the product is a site (type fiber assignment).
  \item Morphisms connect fibers that share at least one type component.
\end{enumerate}

\section{Local Equivalence Check}

For each site $U_i$, the Z3 solver checks:
\[
  \exists \vec{p}.\; \bigwedge_k \mathsf{tag}(p_k) = \tau_k^{(i)}
  \;\land\; \mathsf{tag}(p_k) \neq \mathsf{TBot}
  \;\land\; \den{f}(\vec{p}) \neq \den{g}(\vec{p})
\]

If UNSAT: local equivalence holds ($\Iso(U_i) = 1$).
If SAT: counterexample found ($\Iso(U_i) = 0$).

The solver includes all semantic axioms (Section~\ref{ch:axioms}),
enabling Z3 to use the path constructors during proof search.

\section{\v{C}ech Cohomology Computation}

The local judgments feed into the \texttt{CechCohomologyComputer}:
\begin{enumerate}
  \item Build $C^0$ from local judgments.
  \item Compute $\delta^0$ via \texttt{CoboundaryOperator.d0}.
  \item Compute $\delta^1$ via \texttt{CoboundaryOperator.d1}.
  \item Compute $\Hone = \ker(\delta^1) / \mathrm{im}(\delta^0)$
    via GF(2) rank computation.
  \item If $\Hone = 0$: EQUIVALENT.
  \item If $\Hone \neq 0$: INEQUIVALENT with obstruction data.
\end{enumerate}


